\section{Maximum Principles}\label{chap:maximum-principles}

The maximum principle asserts that solutions of certain scalar elliptic equations of second order cannot have a maximum (or a minimum) in the interior of the domain where they are defined. The basic idea is quite simple. Consider, for simplicity, Laplace\textquotesingle s equation $\Delta u=0$. If $u$ has a maximum at a point $\bm{x}$ and the second derivatives of $u$ do not all vanish at $\bm{x}$, then $\Delta u$ is negative at $\bm{x}$, in contradiction to the equation. The only case left to be ruled out is that of degenerate maxima where all second derivatives vanish. This is accomplished by an approximation argument which removes the degeneracy.

The maximum principle can be used to show that solutions of certain equations must be non-negative. This is important for quantities which have a physical interpretation as densities, concentrations, probabilities, etc. The maximum principle also leads to easy uniqueness results. In later chapters we shall see that in certain problems uniqueness also implies existence. The maximum principle itself can also be used to construct existence proofs. In the next section, we shall give Perron\textquotesingle s existence proof for Dirichlet\textquotesingle s problem. A very recent application of the maximum principle, too complicated to be discussed here, concerns "viscosity solutions" for Hamilton-Jacobi equations. In the third section of this chapter, we shall discuss a result of Gidas, Ni and Nirenberg~\cite{GNN}, which asserts that positive solutions to certain elliptic boundary-value problems must be radially symmetric. The final section of the chapter is concerned with the extension of the maximum principle to parabolic equations.

\subsection{Maximum Principles of Elliptic Problems}\label{sec:elliptic-maximum-principles}

\subsubsection{The Weak Maximum Principle}\label{subsec:weak-maximum-principle}

Throughout this section, we shall consider a second-order operator of the form
\begin{equation}\label{eq:elliptic-second-order-operator}
Lu=a_{ij}(\bm{x})\dfrac{\partial^2u}{\partial x_i\partial x_j}+b_i(\bm{x})\dfrac{\partial u}{\partial x_i}+c(\bm{x})u.
\end{equation}


The following assumptions are made throughout and will therefore not be stated with each theorem. $\Omega$ is a domain in $\mathbb{R}^n$. The coefficients $a_{ij}$, $b_i$ and $c$ are continuous on $\overline\Omega$, and $u$ is in $C^2(\Omega)\cap C(\overline\Omega)$. The matrix $a_{ij}$ is symmetric and strictly positive definite at every point $\bm{x}\in\overline\Omega$, i.e., $L$ is elliptic.

The weak maximum principle is expressed by the following theorem.

\begin{Theorem}
\label{thm:weak-maximum-principle}
Assume that $Lu\geq0$ (or, respectively, $Lu\leq0$) in a bounded domain $\Omega$ and that $c(\bm{x})=0$ in $\Omega$. Then the maximum (or, respectively, the minimum) of $u$ is achieved on $\partial\Omega$.
\end{Theorem}

\begin{Proof}
If $Lu>0$ in $\Omega$, then $u$ cannot achieve its maximum anywhere in $\Omega$. Suppose it did, say at the point $\bm{x}_0$. Then all first derivatives of $u$ vanish at this point, and hence
\begin{equation}\label{eq:elliptic-operator-at-critical-point}
Lu=a_{ij}\dfrac{\partial^2u}{\partial x_i\partial x_j}.
\end{equation}
But at a maximum the matrix of second partial derivatives is negative semidefinite and we conclude (see Problem~\ref{ex:trace-positive-semidefinite-product}) that $Lu(\bm{x}_0)\leq0$, a contradiction.

For the general case, consider the function $u_\epsilon=u+\epsilon\exp(\gamma x_1)$. We find
\begin{equation}\label{eq:weak-maximum-exponential-perturbation}
Lu_\epsilon=Lu+\epsilon(\gamma^2a_{11}+\gamma b_1)\exp(\gamma x_1).
\end{equation}
We now choose $\gamma$ large enough so that $\gamma^2a_{11}+\gamma b_1>0$ throughout $\Omega$ (this is possible since $a_{11}$ is positive and continuous on $\overline\Omega$). Then $Lu_\epsilon>0$ for any positive $\epsilon$. We conclude that
\begin{equation}\label{eq:weak-maximum-perturbed-boundary-equality}
\max_{\overline\Omega}u_\epsilon=\max_{\partial\Omega}u_\epsilon.
\end{equation}
The theorem follows by letting $\epsilon\to0$.
\end{Proof}

\begin{Remark}
\label{rem:weak-maximum-semidefinite}
For later use in connection with parabolic equations, we remark that the proof of Theorem~\ref{thm:weak-maximum-principle} still works if the matrix $a_{ij}$ is only positive semidefinite, as long as there is at least one vector $\bm{\xi}$ independent of $\bm{x}\in\overline\Omega$ such that $\xi_i a_{ij}\xi_j>0$.
\end{Remark}

We have the following corollary of Theorem~\ref{thm:weak-maximum-principle}.

\begin{Corollary}
\label{cor:weak-maximum-nonpositive-zero-order}
Let $\Omega$ be bounded and assume $c\leq0$ in $\Omega$. Let $Lu\geq0$ (or, respectively, $Lu\leq0$). Then
\begin{equation}\label{eq:weak-maximum-positive-negative-parts}
\max_{\overline\Omega}u\leq\max_{\partial\Omega}u^+\quad\text{(or, resp., }\min_{\overline\Omega}u\geq\min_{\partial\Omega}u^-\text{)}.
\end{equation}
Here, $u^+=\max(u,0)$, $u^-=\min(u,0)$. In particular, if $Lu=0$ in $\Omega$, then
\begin{equation}\label{eq:weak-maximum-absolute-value-boundary}
\max_{\overline\Omega}|u|=\max_{\partial\Omega}|u|.
\end{equation}

\end{Corollary}

\begin{Proof}
If $u\leq0$ throughout $\Omega$, the corollary is trivially true. Hence we may assume that $\Omega^+=\Omega\cap\{u>0\}\ne\emptyset$. On $\Omega^+$, we have $-cu\geq0$, and hence
\begin{equation}\label{eq:positive-set-zero-order-free-inequality}
a_{ij}\dfrac{\partial^2u}{\partial x_i\partial x_j}+b_i\dfrac{\partial u}{\partial x_i}\geq0.
\end{equation}
Hence the previous theorem implies that the maximum of $u$ on the closure of $\Omega^+$ is equal to its maximum on $\partial\Omega^+$. Since $u=0$ on $\partial\Omega^+\cap\Omega$, this maximum must be achieved on $\partial\Omega$.
\end{Proof}

The following corollary is typically used in applications. It yields a uniqueness result as well as a comparison principle.

\begin{Corollary}
\label{cor:elliptic-comparison}
Let $\Omega$ be bounded and $c\leq0$. If $Lu=Lv$ in $\Omega$ and $u=v$ on $\partial\Omega$, then $u=v$ in $\Omega$. If $Lu\leq Lv$ in $\Omega$ and $u\geq v$ on $\partial\Omega$, then $u\geq v$ in $\Omega$.
\end{Corollary}

\begin{Remark}
\label{rem:elliptic-comparison-sturm-liouville}
We draw the reader\textquotesingle s attention to the particular case $v=0$. The reader should also note the relationship between this result and the oscillation and comparison theorems of Sturm-Liouville theory in ODEs (cf.~\cite{In}).
\end{Remark}

We conclude this subsection with a definition.

\begin{Definition}
\label{def:elliptic-subsolution}
Assume that $c\leq0$. If $Lu\geq0$ (or, resp., $Lu\leq0$), then $u$ is called a \textbf{subsolution} (\textbf{supersolution}) of the equation $Lu=0$. Subsolutions of $\Delta u=0$ are called \textbf{subharmonic}, and supersolutions are called \textbf{superharmonic}.
\end{Definition}

The terminology is motivated by Corollary~\ref{cor:elliptic-comparison}. A subsolution is less than or equal to a solution with the same values on the boundary; a supersolution is greater than or equal to a solution.
